\chapter{Geometric limits of generalized Ricci flows}
\label{chap:geometric-limits}

Throughout, fix $\epsilon>0$ sufficiently small for the Appendix with
$2\epsilon$ and $\alpha=10^{-2}$, and so that Proposition~\ref{narrows}
holds for $2\epsilon$.

\section{A smooth blow-up limit defined for a small time}

\begin{theorem}[Existence of a short-time blow-up limit]
\label{thm:short-time-blowup-limit}
\label{smlmtflow}
Fix canonical-neighborhood constants $(C,\epsilon)$ and non-collapsing
constants $r>0,\kappa>0$. Let $({\mathcal M}_n,G_n,x_n)$ be based
generalized $3$-dimensional Ricci flows, put $t_n={\bf t}(x_n)$ and
$Q_n=R(x_n)$, and let $M_n$ be the $t_n$-slice. Assume:
\begin{enumerate}
\item[(1)] Each flow is defined in a time interval contained in
$[0,\infty)$ and has curvature pinched toward positive, or has
non-negative curvature.
\item[(2)] If ${\bf t}(y_n)\le t_n$ and $R(y_n)\ge4R(x_n)$, then
$y_n$ has a strong $(C,\epsilon)$-canonical neighborhood.
\item[(3)] $\lim_{n\to\infty}Q_n=\infty$.
\item[(4)] For every $A<\infty$, eventually
$B(x_n,t_n,AQ_n^{-1/2})$ has compact closure in $M_n$ and the flow is
$\kappa$-non-collapsed there on scales at most $r$.
\item[(5)] There is $\mu>0$ such that, for every $A<\infty$, eventually
every $y_n\in B(x_n,t_n,AQ_n^{-1/2})$ has a backward flow line of length at
least $\mu(\max(Q_n,R(y_n)))^{-1}$.
\end{enumerate}
After a subsequence and shifting $t_n$ to $0$, the based manifolds
$(M_n,Q_nG_n(0),x_n)$ converge geometrically to a complete $3$-manifold
$(M_\infty,g_\infty,x_\infty)$ of bounded non-negative curvature. For some
$t_0>0$, depending on the curvature bound of the limit and on $\mu$, the
flows converge to a limit flow on $[-t_0,0]$ with final metric
$g_\infty(0)=g_\infty$.
\end{theorem}

\begin{proof}
\uses{bcbd,2ndmfdconv,blowupposcurv,partialflowlimit}
\sketch Shift time so $t_n=0$. The bounded-curvature-at-bounded-distance
theorem gives, for each $A$, a scalar-curvature bound $Q(A)$ on
$B_{Q_nG_n}(x_n,0,A)$. The following local estimate supplies a uniform
backward parabolic bound.
\end{proof}

\begin{lemma}[Curvature control from canonical neighborhoods]
\label{lem:curvature-control-canonical-nbhd}
\label{controlnbhd}
Let $({\mathcal M},G)$ be a generalized $3$-dimensional Ricci flow. Suppose
$r_0>0$ and every point with $R\ge r_0^{-2}$ has a strong
$(C,\epsilon)$-canonical neighborhood. For ${\bf t}(z)=t_0$, set
$$r=\frac{1}{2C\sqrt{\max(R(z),r_0^{-2})}},\qquad
\Delta t=\frac{1}{16C(R(z)+r_0^{-2})}.$$
If $r'\le r$, $|t'-t_0|\le\Delta t$, and
$j:B(z,t_0,r')\times I\to{\mathcal M}$ is a time- and vector-field-compatible
embedding, where $I$ has endpoints $t_0,t'$, then
$$R(y)\le2(R(z)+r_0^{-2})$$
on the image of $j$.
\end{lemma}

\begin{proof}
  \uses{bcbd, 2ndmfdconv, blowupposcurv, partialflowlimit, lem:curvature-control-canonical-nbhd}
\sketch Along a spatial arclength path, the canonical-neighborhood estimate
$|R'|\le CR^{3/2}$ gives
\begin{equation}\label{Rineq}
R(y)\le\frac{16}{9}(R(z)+r_0^{-2}).
\end{equation}
Along a vertical flow line, $|R'|\le CR^2$; the three cases according to
the closed set where $R\ge r_0^{-2}$ give respectively the bounds
$r_0^{-2}$, $16r_0^{-2}/15$, and $2(R(z)+r_0^{-2})$, using the lengths
$r_0^2/(16C)$ and $\Delta t$.
\end{proof}

\begin{corollary}[Uniform curvature bound on a backward parabolic neighborhood]
\label{cor:uniform-backward-curvature-bound}
\uses{lem:curvature-control-canonical-nbhd}
\label{cnnbhd2}

For each $A<\infty$, if $Q(A)$ bounds the scalar curvature of
$Q_nG_n$ on $B_{Q_nG_n}(x_n,0,A)$, then there are $t'_0(A)>0$ and
$Q'(A)<\infty$, with $Q'(A)$ depending only on $Q(A)$, such that eventually
$B_{Q_nG_n}(x_n,0,A)\times(-t'_0(A),0]$ embeds compatibly in
${\mathcal M}_n$ and has scalar curvature at most $Q'(A)$.
\end{corollary}
\begin{proof}\uses{lem:curvature-control-canonical-nbhd}\sketch Apply Lemma~\ref{lem:curvature-control-canonical-nbhd} and
Assumption (5) of the theorem.\end{proof}

\begin{corollary}[Bounded covariant derivatives of curvature along the blow-up sequence]
\label{cor:shi-bound-blowup-sequence}
\uses{shi, cor:uniform-backward-curvature-bound}

For each $A<\infty$ and integer $\ell\ge0$, there is $C_2$ such that
eventually
$$|\nabla^\ell\Rm_{Q_nG_n}(x)|\le C_2$$
for $x\in B_{Q_nG_n}(x_n,0,A)$.
\end{corollary}
\begin{proof}\sketch Apply Shi's theorem to the preceding parabolic bounds.\end{proof}

\begin{claim}[Non-collapsing volume bound for the blow-up sequence]
\label{claim:volume-noncollapsing-blowup}
There is $\eta>0$ such that eventually
$$\Vol(B_{Q_nG_n}(x_n,0,\eta))\ge\kappa\eta^3.$$
\end{claim}
\begin{proof}\sketch Rescaling preserves the stated $\kappa$-non-collapsing
hypothesis, while the curvature bounds give the required fixed scale.\end{proof}

\begin{claim}[Points of large curvature in the limit have canonical neighborhoods]
\label{claim:canonical-neighborhood-persists}
\uses{thm:short-time-blowup-limit}

Every point of $(M_\infty,g_\infty)$ with curvature greater than $4$ has a
$(2C,2\epsilon)$-canonical neighborhood.
\end{claim}
\begin{proof}\uses{thm:short-time-blowup-limit,canonvary}
\sketch Exhaust the limit by the geometric convergence embeddings. The
corresponding points have $(C,\epsilon)$-canonical neighborhoods of uniformly
bounded diameter, and Proposition~\ref{canonvary} passes them to the limit.\end{proof}

\begin{claim}[The short-time blow-up limit has bounded curvature]
\label{claim:limit-bounded-curvature}
The limit $(M_\infty,g_\infty)$ has bounded curvature.
\end{claim}
\begin{proof}\uses{claim:canonical-neighborhood-persists,localprod,narrows}
\sketch If a curvature eigenvalue vanishes, Corollary~\ref{localprod} gives a
one- or two-sheeted product cover, hence bounded curvature. Compact limits are
bounded. In the remaining noncompact positive-curvature case, unbounded
curvature would produce $2\epsilon$-necks of arbitrarily small scale, contrary
to Proposition~\ref{narrows}.\end{proof}

\section{Long-time blow-up limits}

\begin{theorem}[Existence of a long-time (ancient) blow-up limit]
\label{thm:long-time-blowup-limit}
\uses{thm:short-time-blowup-limit}
\label{kaplimit}

Suppose the hypotheses of Theorem~\ref{thm:short-time-blowup-limit} hold.
Suppose also that $0<T_0\le\infty$ and, for every $T<T_0$ and $A<\infty$,
eventually $B(x_n,t_n,AQ_n^{-1/2})\times(t_n-TQ_n^{-1},t_n]$ embeds compatibly
and is $\kappa$-non-collapsed on scales at most $r$. Then a subsequence of
$(Q_n{\mathcal M}_n,Q_nG_n,x_n)$ converges, after $t_n=0$, to a complete
non-negative-curvature flow on $(-T_0,0]$, locally bounded in time. If
$T_0=\infty$, the limit is a $\kappa$-solution.
\end{theorem}
\begin{remark}[Comparison with the short-time blow-up limit]
\label{rem:comparison-short-and-long-time-limits}
\uses{thm:short-time-blowup-limit, thm:long-time-blowup-limit}

The short-time interval depends on the final-slice curvature, whereas the
interval here depends only on the available backward time.
\end{remark}
\begin{proof}
  \uses{thm:short-time-blowup-limit, blowupposcurv, narrows, Harnack, thm:long-time-blowup-limit, lem:curvature-control-canonical-nbhd, cor:curvature-bound-near-basepoint, claim:distance-distortion-bound, bcbd, cor:distance-integral-bound}
\sketch Start with Theorem~\ref{thm:short-time-blowup-limit} and repeatedly
extend the limit backward using the following proposition. Harnack controls
the resulting time-slices. When $T_0=\infty$, Proposition~\ref{narrows} and
Harnack give bounded curvature, the rescaled non-collapsing scales tend to
infinity, and $R(x_\infty,0)=1$, so the limit is a nonflat $\kappa$-solution.
\end{proof}

\begin{proposition}[Extending a bounded-curvature limit flow backward in time]
\label{prop:extend-limit-flow-backward}
\uses{thm:long-time-blowup-limit}
\label{maxT}

Under the hypotheses of Theorem~\ref{thm:long-time-blowup-limit}, if a complete
non-negative-curvature limit flow is defined on $(-T,0]$, $T<T_0$, and is
locally bounded in time, then its curvature is bounded and it extends with
bounded curvature to $(-(T+\delta),0]$ for some $\delta>0$.
\end{proposition}
\begin{proof}\uses{lem:curvature-control-canonical-nbhd,thm:short-time-blowup-limit,cor:curvature-bound-near-basepoint,narrows,Harnack}
\sketch If the scalar curvature is bounded, Lemma~\ref{lem:curvature-control-canonical-nbhd}
and Proposition~\ref{partialflowlimit} extend the flow uniformly beyond $-T$.
The required bound follows from the compact and noncompact alternatives in
the next lemma and its claims, followed by Harnack.
\end{proof}

\begin{lemma}[Curvature bound on a compact subset from a pointwise minimum bound]
\label{lem:curvature-bound-compact-subset}\label{compactbd}
Suppose a complete non-negative-curvature limit flow is defined on $(-T,0]$,
$T<\infty$, with curvature locally bounded in time. If $X$ is compact and
connected and $\min_XR_{g_\infty}(\cdot,t)$ is uniformly bounded, then $R$ is
uniformly bounded on $X\times(-T,0]$.
\end{lemma}
\begin{proof}\uses{claim:distance-distortion-bound,bcbd}\sketch Use the
distance estimate below to obtain a uniform diameter bound, then lift points
of bounded curvature and points of unbounded curvature to the approximating
flows. The bounded-distance curvature theorem gives a contradiction.\end{proof}

\begin{claim}[Distance distortion bound under the limit flow]
\label{claim:distance-distortion-bound}\label{distch}
Let $Q$ bound $R(x,0)$ on $M_\infty$. For $x,y\in M_\infty$ and
$t\in(-T,0]$,
$$d_t(x,y)\le d_0(x,y)+16\sqrt{\frac Q3}T.$$
\end{claim}
\begin{proof}\uses{cor:distance-integral-bound,Harnack}\sketch Harnack gives
$$\frac{\partial R}{\partial t}\ge-\frac R{t+T},\qquad
R(x,-t_0)\le Q\frac T{T-t_0},$$
and hence, in dimension $3$,
$$\Ric(x,-t_0)\le(n-1)\frac{QT}{2(T-t_0)}.$$
Integrating Corollary~\ref{cor:distance-integral-bound} yields
$$d_{-t_0}(x,y)\le d_0(x,y)+8\int_{-t_0}^0\sqrt{\frac{QT}{3(T+t)}}
\le d_0(x,y)+16\sqrt{\frac Q3}T.$$
\end{proof}

\begin{lemma}[Existence of a point far from a distant point in non-negative curvature]
\label{lem:far-point-with-distant-antipode}
For a complete connected noncompact non-negative-curvature manifold and
$x_0\in M$, there is $D>0$ such that $d(x_0,y)=d\ge D$ implies the existence
of $x$ with $d(y,x)=d$ and $d(x_0,x)>3d/2$.
\end{lemma}
\begin{proof}\uses{cor:angle-monotonicity}\sketch Otherwise geodesics to a
sequence $y_n$ with $d_n\to\infty$ converge to a ray. The angle comparison
gives $\alpha_n/d_n\to0$, and a point $z_n$ on the ray satisfies
$d(y_n,z_n)=d_n$ and $d(x_0,z_n)\ge2d_n-\alpha_n>3d_n/2$, a contradiction.\end{proof}

\begin{claim}[Curvature bound outside a large ball around the base point]
\label{claim:curvature-bound-outside-large-ball}\label{outerbd}
\uses{lem:far-point-with-distant-antipode}

If $D$ is at least the constant in the preceding lemma and
$D\ge32\sqrt{Q/3}T$, then $R_{g_\infty}(y,t)$ is uniformly bounded for
$y\notin B(x_\infty,0,D)$ and $t\in(-T,0]$.
\end{claim}
\begin{proof}\uses{claim:distance-distortion-bound,lem:far-point-with-distant-antipode}
\sketch Choose $z$ with $d_0(y,z)=d$ and $d_0(x_\infty,z)>3d/2$. An
unbounded curvature sequence at $y$ gives a small canonical neighborhood.
If the two minimizing geodesics exit the same end, Toponogov and
$$\cos\theta\ge1-\frac{(4\pi\epsilon)^2}{2},\qquad
d_t(x_\infty,y),d_t(y,z)\le d+16\sqrt{Q/3}T$$
contradict $d_0(y,z)=d$ because $\epsilon<1/(8\pi)$. The neck alternative
then separates $x_\infty$ and $z$ by balls of radius
$4\pi R(y,t)^{-1/2}$; arbitrarily small such balls are impossible.\end{proof}

\begin{claim}[The point is a neck center with geodesics exiting opposite ends]
\label{claim:neck-center-opposite-ends}
\uses{claim:distance-distortion-bound}

The point $y$ in the preceding claim cannot lie in a cap; it is the center of
a $2\epsilon$-neck, and minimizing geodesics from $y$ to $x_\infty$ and $z$
exit opposite ends.
\end{claim}
\begin{proof}\sketch A canonical neighborhood has diameter at most
$2CR(y)^{-1/2}$. Equal-radius points on geodesics exiting the same end satisfy
$d_t(b,y)/d_t(a,y)<4\pi\epsilon$, so Toponogov gives the contradiction used
above.\end{proof}

\begin{corollary}[Curvature bound near the base point for all backward times]
\label{cor:curvature-bound-near-basepoint}\label{ancientlimit}
\uses{lem:curvature-bound-compact-subset, claim:curvature-bound-outside-large-ball}

Any complete non-negative-curvature limit flow on $(-T,0]$, $T<\infty$,
locally bounded in time, has uniformly bounded scalar curvature on
$B(x_\infty,0,A)\times(-T,0]$ for every $A<\infty$.
\end{corollary}
\begin{proof}\uses{lem:curvature-bound-compact-subset,claim:curvature-bound-outside-large-ball,scalarevol}
\sketch If the limit is compact, $\min R$ is controlled by
Proposition~\ref{scalarevol} and Lemma~\ref{lem:curvature-bound-compact-subset}.
If noncompact, first use the preceding exterior bound on the closure of a
large ball, then apply the compact-subset lemma to each larger ball.\end{proof}

\begin{claim}[Uniform curvature bound on an extended backward time interval]
\label{claim:uniform-curvature-bound-extended-interval}
\uses{cor:curvature-bound-near-basepoint}

For each $A<\infty$, eventually there are $\delta>0$, $\delta\le T_0-T$, and
a bound independent of $n$ for the scalar curvature of
$B_{Q_nG_n}(x_n,0,A)\times[-(T+\delta),0]$.
\end{claim}
\begin{proof}\uses{lem:curvature-control-canonical-nbhd,cor:curvature-bound-near-basepoint}
\sketch Apply Corollary~\ref{cor:curvature-bound-near-basepoint} on the
$2A$-ball and Lemma~\ref{lem:curvature-control-canonical-nbhd}; then use
\ref{narrows}, Hamilton's compactness theorem, and Harnack to extend the
limit and bound its curvature.\end{proof}

\section{Incomplete smooth limits at singular times}
\subsection{Assumptions}
\begin{definition}[Standing assumptions for limits at a singular time]
\label{def:standing-assumptions-singular-time}\label{assumptions}
\begin{enumerate}
\item[(a)] Singular times are discrete in $\Ar$, and nonsingular time-slices
are compact $3$-manifolds.
\item[(b)] The flow is defined in $[0,\infty)$ and curvature is pinched
toward positive.
\item[(c)] For some $r_0>0,C<\infty$, points with $R\ge r_0^{-2}$ have
strong $(C,\epsilon)$-canonical neighborhoods, so
$$\left|\frac{\partial R}{\partial t}\right|<CR^2,\qquad
|\nabla R|<CR^{3/2}.$$
\end{enumerate}
\end{definition}

\begin{theorem}[Properties of the limit metric at a singular time]
\label{thm:singular-time-limit-properties}\label{omega}
\uses{def:standing-assumptions-singular-time}

Let $({\mathcal M},G)$ satisfy the assumptions above on $0\le t<T<\infty$.
Choose $T^-<T$, identify the time interval with $M=M_{T^-}$, and write
$g(t)$ for the pulled-back metrics. Set
$$\Omega=\{x\in M\mid\liminf_{t\to T}R_g(x,t)<\infty\}.$$
Then $\Omega$ is open, $g(t)|_\Omega$ converges smoothly on compact sets to
$g(T)$, $R(g(T)):\Omega\to\Ar$ is proper and bounded below, and the maximal
extension
$$\widehat{\mathcal M}={\mathcal M}\cup_{\Omega\times[T^-,T)}
(\Omega\times[T^-,T])$$
is a generalized Ricci flow. Every end of a component of $\Omega$ lies in a
strong $2\epsilon$-tube, and every $(x,T)$ with $R(x,T)>r_0^{-2}$ has a
strong $(2C,2\epsilon)$-canonical neighborhood in $\widehat{\mathcal M}$.
\end{theorem}
\begin{remark}[Proper functions]\label{rem:proper-function}
A function is proper when inverse images of compact sets are compact.
\end{remark}

\begin{lemma}[Omega is open and the metric converges to a limit at time T]
\label{lem:omega-open-limit-metric}\label{Rproper}
\uses{def:standing-assumptions-singular-time, thm:singular-time-limit-properties}

The set $\Omega$ is open, $g(t)|_\Omega$ converges smoothly on compact sets
to $g(T)$, and $R(g(T))$ is proper and bounded below.
\end{lemma}
\begin{proof}\uses{lem:curvature-control-canonical-nbhd,shi}\sketch A bounded
curvature sequence at $x\in\Omega$ and Lemma~\ref{lem:curvature-control-canonical-nbhd}
give a neighborhood with bounded curvature up to $T$. Shi estimates yield
$C^\infty$ convergence. If a sequence approaches $M\setminus\Omega$ while
$R(g(T))$ stays bounded, the same lemma bounds curvature at the limit point,
a contradiction.\end{proof}

\begin{definition}[Maximal extension of a generalized Ricci flow to time T]
\label{def:maximal-extension-to-time-T}\label{hatM}
\uses{thm:singular-time-limit-properties, lem:omega-open-limit-metric}

Set
$$\widehat {\mathcal M}= {\mathcal M}\cup_{\Omega\times[T^-,T)}
(\Omega\times[T^-,T]).$$
The compatible space-time structures and metrics glue to $\widehat G$,
which satisfies Ricci flow and extends $({\mathcal M},G)$; the time-slice
$\Omega$ need not be complete.
\end{definition}

\subsection{Canonical neighborhoods for $(\widehat{\mathcal M},\widehat G)$}
\begin{lemma}[Canonical neighborhoods for points of Omega at time T]
\label{lem:canonical-neighborhoods-at-time-T}\label{omegacanon}
\uses{def:maximal-extension-to-time-T, def:standing-assumptions-singular-time}

If $(x,T)\in\Omega\times\{T\}$ and $R(x,T)>r_0^{-2}$, then exactly one of
the available alternatives holds: $(x,T)$ is the center of a strong
$2\epsilon$-neck; it lies in a $(2C,2\epsilon)$-cap core; it lies in a
$2C$-component; or it lies in a $2\epsilon$-round component.
\end{lemma}
\begin{proof}\uses{lem:curvature-control-canonical-nbhd,lem:omega-open-limit-metric,canonvary}
\sketch Canonical neighborhoods at times tending to $T$ lie in a fixed
compact subset and their metrics converge smoothly. Caps and round or
$C$-components persist with doubled parameters. The neck case is the claim
below.\end{proof}

\begin{claim}[A limit of strong epsilon-necks is a strong 2epsilon-neck]
\label{claim:limit-of-necks-is-neck}\label{epslimit}
If $x\in\Omega$ is the center of strong $\epsilon$-necks at times
$t_n\to T$, then $(x,T)$ is the center of a strong $2\epsilon$-neck.
\end{claim}
\begin{proof}\sketch Rescale so $R(x,T)=1$. On
$B(x,T,2\epsilon^{-1}/3)$ the neck flow lines exist on $(T-1,T]$ and the
rescaled neck metrics converge smoothly. Restricting the neck coordinates to
$S^2\times(-\epsilon^{-1}/2,\epsilon^{-1}/2)$ gives a strong
$2\epsilon$-neck at time $T$.\end{proof}

\subsection{The ends of $(\Omega,g(T))$}
\begin{definition}[Strong 2ε-horn]
\label{def:strong-epsilon-horn}
\uses{def:maximal-extension-to-time-T}
A strong $2\epsilon$-horn is a proper copy of $S^2\times[0,1)$ in $\Omega$
whose points are centers of strong $2\epsilon$-necks and whose boundary is a
central neck sphere.
\end{definition}
\begin{definition}[Strong double 2ε-horn and capped 2ε-horn]
\label{def:strong-double-and-capped-horn}
\uses{def:strong-epsilon-horn}
A strong double $2\epsilon$-horn is a component diffeomorphic to
$S^2\times(0,1)$ consisting entirely of strong $2\epsilon$-necks, hence a
$2\epsilon$-tube with two horn ends. A $C'$-capped $2\epsilon$-horn is the
union of a $(C',2\epsilon)$-cap core and such a horn; it is an open
$3$-ball or a punctured $\Ar P^3$.
\end{definition}
\begin{definition}[The subset $\Omega_\rho$ and horns with boundary contained in it]
\label{def:omega-rho}
\uses{thm:singular-time-limit-properties, def:strong-epsilon-horn}
For $0<\rho<r_0$, let
$$\Omega_\rho=\{x\in\Omega\mid R(x,T)\le\rho^{-2}\}.$$
A horn has boundary contained in $\Omega_\rho$ when its boundary sphere is
contained in this set.
\end{definition}

\begin{lemma}[Classification of components of Omega disjoint from $\Omega_\rho$]
\label{lem:omega-component-classification}
\uses{lem:canonical-neighborhoods-at-time-T, def:omega-rho}
If a component $\Omega^0$ contains no point of $\Omega_\rho$, it is one of:
\begin{enumerate}
\item a strong double $2\epsilon$-horn, diffeomorphic to $S^2\times\Ar$;
\item a $2C$-capped horn, diffeomorphic to $\Ar^3$ or punctured $\Ar P^3$;
\item a compact union of two cap cores and a tube, diffeomorphic to
$S^3$, $\Ar P^3$, or $\Ar P^3\#\Ar P^3$;
\item a compact $2\epsilon$-round component of constant positive curvature;
\item a compact component fibering over $S^1$ with fiber $S^2$; or
\item a compact $2C$-component diffeomorphic to $S^3$ or $\Ar P^3$.
\end{enumerate}
\end{lemma}
\begin{proof}\uses{lem:canonical-neighborhoods-at-time-T,epstopology}\sketch Every
point has a $(2C,2\epsilon)$-canonical neighborhood. The round and component
alternatives give cases (4) and (6); otherwise Proposition~\ref{epstopology}
classifies the neck and cap alternatives as (1)--(3) or (5).\end{proof}
\begin{remark}[Possibly infinitely many double horns]
\label{rem:possibly-infinitely-many-horns}
\uses{def:strong-double-and-capped-horn}
The hypotheses do not assert that the number of double $2\epsilon$-horns is finite.
\end{remark}

\begin{lemma}[Finitely many horns meet $\Omega_\rho$]
\label{lem:finitely-many-horns-meeting-omega-rho}\label{2epshorns}
\uses{def:omega-rho, def:standing-assumptions-singular-time}

For a flow satisfying the standing assumptions and $0<\rho<r_0$, the union
$\Omega^0(\rho)$ of components meeting $\Omega_\rho$ has finitely many
components and is a compact set together with finitely many strong
$2\epsilon$-horns, disjoint from $\Omega_\rho$, whose boundaries lie in
$\Omega_{\rho/2C}$.
\end{lemma}
\begin{proof}\uses{lem:canonical-neighborhoods-at-time-T,lem:omega-component-classification,def:omega-rho,Xcontainedin}
\sketch Properness makes $\Omega_\rho$ compact. In a noncompact component,
the component of $\{R\ge2C^2\rho^{-2}\}$ containing an end is, by
Proposition~\ref{Xcontainedin}, a strong tube rather than a capped tube.
Its frontier neck has curvature $2C^2\rho^{-2}$ and central sphere in
$\Omega_{\rho/2C}$. Distinct end horns are disjoint unless they force a
two-ended $S^2\times(0,1)$ component; compactness of
$\Omega_{\rho/2C}$ then gives finiteness.\end{proof}

\section{Existence of strong $\delta$-necks sufficiently deep in a $2\epsilon$-horn}
\begin{theorem}[Strong $\delta$-necks sufficiently deep in a 2ε-horn]
\label{thm:deep-delta-necks-in-horn}\label{hexist}
\uses{def:omega-rho, def:standing-assumptions-singular-time}

For fixed $\rho>0$ and every $\delta>0$, there is
$0<h=h(\delta,\rho)\le\min(\rho\delta,\rho/2C)$, depending implicitly on
$r,(C,\epsilon)$, such that every point $x$ in a $2\epsilon$-horn with
boundary in $\Omega_{\rho/2C}$ and $R(x,T)\ge h^{-2}$ centers a strong
$\delta$-neck contained in the horn. There is also $y$ in the horn with
$R(y)=h^{-2}$ whose neck sphere cuts off an end disjoint from $\Omega_\rho$.
\end{theorem}
\begin{proof}
  \uses{thm:short-time-blowup-limit, thm:splitting-theorem, blowupposcurv, localprod, rm>0, 2DGSS, bcbd}
\sketch If no such $h$ exists, choose horns and points with
$Q_n=R(x_n)\to\infty$ lacking strong $\delta$-necks. The rescaled sequence
satisfies the five hypotheses below, so it has a nonflat non-negative-curvature
limit. The splitting claim makes the limit $(S^2,h)\times(\Ar,ds^2)$;
backward extension makes it ancient, and Corollary~\ref{2DGSS} makes it the
shrinking round cylinder, contradicting the choice of $x_n$. The level set
$R=h^{-2}$ and $h<\rho/2C$ give the final separating neck.\end{proof}

\begin{claim}[The rescaled sequence satisfies the short-time blow-up hypotheses]
\label{claim:rescaled-sequence-satisfies-hypotheses}\label{hyp}
\uses{thm:short-time-blowup-limit}

The rescaled sequence in the proof of Theorem~\ref{thm:deep-delta-necks-in-horn}
satisfies all five hypotheses of Theorem~\ref{thm:short-time-blowup-limit}.
\end{claim}
\begin{proof}\uses{thm:short-time-blowup-limit}\sketch Hypotheses (1) and
(3) are inherited. The canonical-neighborhood threshold gives (2), and horn
geometry gives non-collapsing and backward flow lines for (4)--(5).\end{proof}

\begin{claim}[Balls around $x_n$ eventually lie in the horn with compact closure]
\label{claim:ball-contained-in-horn-compact-closure}
For every $A<\infty$, eventually
$B(x_n,0,AQ_n^{-1/2})$ lies in the $2\epsilon$-horn and has compact closure.
\end{claim}
\begin{proof}\uses{bcbd}\sketch Boundary points have curvature at most
$16C^2\rho^{-2}$ and their necks control curvature by $32C^2\rho^{-2}$ within
distance $\epsilon^{-1}\rho/2C$. Hence the rescaled distance to the boundary
tends to infinity. Distance to the horn end also tends to infinity, since
otherwise Theorem~\ref{bcbd} is violated.\end{proof}

\begin{claim}[The blow-up limit splits as a round $S^2$ times $\mathbb{R}$]
\label{claim:limit-splits-as-cylinder}
$(M_\infty,g_\infty(0))$ is isometric to $(S^2,h)\times(\Ar,ds^2)$, where
$h$ has non-negative curvature.
\end{claim}
\begin{proof}\uses{thm:splitting-theorem}\sketch Rays from $x_n$ to the horn
end and to its boundary have lengths tending to infinity and converge to two
ends of the limit. The splitting theorem gives a product with a non-negative
curvature surface; nonflatness makes that surface diffeomorphic to $S^2$.\end{proof}

\begin{claim}[Backward flow lines with controlled curvature near the cylinder limit]
\label{claim:backward-flow-line-curvature-bound}
For $z\in M_\infty$, eventually $\varphi_n(z)$ has a backward flow line in
$(Q_n{\mathcal M}_n,Q_nG_n)$ on
$$\left(-T-\frac{R_{Q_nG_n}^{-1}(\varphi_n(z),0)}2,0\right]$$
and the scalar curvature there is at most $R(\varphi_n(z),0)$.
\end{claim}
\begin{proof}\uses{neckcurv,defncanonnbhd}\sketch Smooth convergence on compact
product neighborhoods and positivity of the cylinder limit force the
approximating points to be strong $\epsilon$-neck centers, not cap cores.
The neck curvature evolution and its backward extension give the stated
interval and upper bound.\end{proof}

\begin{corollary}[Monotone choice of $h(\rho,\delta)$]
\label{cor:monotone-h-function}\label{hmonotone}
\uses{thm:deep-delta-necks-in-horn}

The function $h(\rho,\delta)$ may be chosen with $h\le\delta\rho$, weakly
non-decreasing in $\delta$ for fixed $\rho$, and weakly non-decreasing in
$\rho$ for fixed $\delta$.
\end{corollary}
\begin{proof}\uses{thm:deep-delta-necks-in-horn}\sketch If $h$ works for
$(\rho,\delta)$, it works for larger parameters and every smaller positive
value. Choose decreasing sequences $\rho_n,\delta_n\to0$ and decreasing
$h_n\le\rho_n\delta_n$, then define $h(\rho,\delta)$ by the largest index
with $\rho\ge\rho_n$ and $\delta\ge\delta_n$.\end{proof}
